%% kapa: exam=1 round=30 year=2019 period=1 form=A subject=economics subject_ko=경제학원론 seq=40 number=80
%% source: Q-Net 감정평가사 기출문제 (cst00309) artlSeq=5205205 file=extracted/1008/5205205/2019년 제30회 감정평가사 1차시험 문제지/2019년 제30회 감정평가사 1차시험 문제지/2019년도 제30회 감정평가사 1차 1교시 A형.hwp
%% conversion: kordoc 4.24.51 -> tools/gen1.py (auto); review=proofread-sonnet; flags=none
\begin{question}{80}
모든 시장이 완전경쟁적인 甲국에서 대표적인 기업 A의 생산함수가 ${ Y=4L^{ 0.5 } K^{ 0.5 } }$이다. 단기적으로 A의 자본량은 1로 고정되어 있다. 생산물 가격이 2이고 명목임금이 4일 경우, 이윤을 극대화하는 A의 단기 생산량은? (단, ${ Y }$는 생산량, ${ L }$은 노동량, ${ K }$는 자본량이며, 모든 생산물은 동일한 상품이다.)
\begin{choices}
  \item 1
  \item 2
  \item 4
  \item 8
  \item 16
\end{choices}
\end{question}
\answer{3}
